\section{Overflow indicators and marker shapes}
\key{overflow mark=none|arrow|slash} indicates clipped values at scale boundaries.
The default is none for individual graphics and arrow in \key{gradtable}. Exact boundary
values do not produce a mark. \key{overflow=error} still rejects out-of-range data.
\begin{demo}{Clipping without hiding the original value}
\gradbar[max=100,width=65mm,precision=0,overflow mark=arrow]{125}

\gradrange[min=-40,max=60,width=65mm,precision=0,
  overflow mark=slash]{-55}{75}
\end{demo}
Checks use raw values, interval and error endpoints, both comparison values and
separate signed stack subtotals. Fixed-range sparklines check all present points.
The mark indicates overflow, not quality or the amount beyond the boundary.

\key{marker shape} supports circle, square, diamond and triangle;
\key{marker fill} supports solid and hollow. Dumbbell reference markers and
interval estimates remain hollow to preserve their meaning. Marker size is the
distance from the center to its boundary.
\begin{demo}{Shapes that remain distinguishable in monochrome}
\gradbar[shape=lollipop,marker shape=square]{72}

\graddumbbell[marker shape=diamond]{48}{72}

\gradbar[shape=lollipop,marker shape=triangle,marker fill=hollow]{72}
\end{demo}
